\section{Conclusion}

We benchmarked ResNet-50, EfficientNet-B0, and ViT-S/16 on HAM10000 under three pre-declared conditions. Focal loss with class-balanced sampling raised ViT-S/16 sensitivity by 8.2 points and lowered the two CNNs, which reversed the point-estimate ranking at 95\% specificity. The paired within-architecture intervals include zero, but the difference-in-differences between each CNN and the ViT excludes zero. The AUC ranking stayed stable, and focal loss roughly tripled expected calibration error. Calibration and threshold placement therefore help explain how the sensitivity ranking can change while the ranking quality of the scores does not. The direction of the effect still depends on the architecture.

The fairness analysis became more specific once the cutoff and the error type were separated. The dark stratum has a higher false positive rate at $L^*$ cutoffs of 55, 58, and 60, while the dark AUC deficit is present only at the strictest cutoff. Adjusted models show that darker skin is associated with more false positives and lighter skin with more false negatives, and that lesion-skin contrast moves both error types in opposite directions. Fairness in this setting is an error-type shift, not a single sensitivity gap.

The XAI analysis showed that all three architectures place less attribution inside the lesion on darker images. False negatives have lower attribution mass than false positives, which explains why the pooled mass-error correlation is weak. HIBA validation reproduced the CNN-transformer gap but showed a 16 to 18 point sensitivity drop, and HIBA contains no Fitzpatrick IV melanomas and no Fitzpatrick V-VI images.

The next steps are a clinical dark-skin evaluation with DDI or PASSION, two additional seeds for the six focal and CE cells, and a second attribution method. The clinical dataset is the main requirement for a precise Fitzpatrick V-VI claim. The additional seeds and the second attribution method would tighten the architecture interaction and test whether the error-type pattern is specific to Grad-CAM.
